% Frugal AI seminar, part 7: how to start. Practices mirror FrugalAI/docs/token_frugality.md: things to try and
% measure, not proven savings. CheatSheet-safe: no floats.

\section[Start]{How to Start}

%%%%%%%%%%%%%%%%%%%%%%%%%%%%%%%%%%%%%%%%%%%%%%%%%%%
\begin{frame}[fragile]\frametitle{}
\begin{center}
{\Large How to Start}

\medskip
{\small (Measure first, then trim)}
\end{center}
\end{frame}

%%%%%%%%%%%%%%%%%%%%%%%%%%%%%%%%%%%%%%%%%%%%%%%%%%%
\begin{frame}[fragile]\frametitle{Measure First}
\begin{itemize}
\item Three numbers per model and device: decode tokens per second, time to first token, peak memory of the engine.
\item One warm-up run, then repeats, and report the median and the range, never one run. Two identical runs differed by about 8\% on a desktop PC.
\item Say whether ``thinking'' was on: a reasoning model can spend almost its whole budget on hidden tokens (Qwen3 1.7B spent 126 of 128).
\item Method: FrugalAI benchmarks/METHODOLOGY.md. The command:
\end{itemize}
\begin{lstlisting}
python -m frugalai_harness.bench --backend lmstudio --model <model id> \
       --device <label> --quantisation Q4_K_M
# one warm-up run is discarded; the CSV holds every run, the host and the timing source
\end{lstlisting}
\end{frame}

%%%%%%%%%%%%%%%%%%%%%%%%%%%%%%%%%%%%%%%%%%%%%%%%%%%
\begin{frame}[fragile]\frametitle{Token-Frugal Practices}
\begin{itemize}
\item \textbf{Lean context}: send only what the step needs; retrieve instead of re-sending history.
\item \textbf{Structured, short prompts}: compact fields instead of chatty prose. Test the saving, do not assume it.
\item \textbf{Constrained answers}: ask for a label or a field where the job allows. Check accuracy: shortening reasoning can hurt hard tasks.
\item \textbf{Step budgets}: cap model calls per task and stop when the state repeats.
\item \textbf{Small specialised models} for each step, and code for what code does well.
\item \textbf{Measure energy}, not only tokens: a wall or USB power meter on your device is the most direct method. Whether energy tracks tokens on a given device is something to test, not assume.
\end{itemize}
{\tiny (Ref: FrugalAI docs/token\_frugality.md; practices to try, not proven savings)}
\end{frame}

%%%%%%%%%%%%%%%%%%%%%%%%%%%%%%%%%%%%%%%%%%%%%%%%%%%
\begin{frame}[fragile]\frametitle{A Ten-Minute Start}
\begin{enumerate}
\item Install LM Studio (free) and download a small model, for example \texttt{gemma-2-2b-it} (Q4\_K\_M, 1.71 GB).
\item Start its local server from the Developer tab.
\item Run the harness benchmark above and keep the CSV.
\item Point the offline wiki recipe at a folder of your own notes and ask it three questions.
\item Run any LangGraph example from the seminar code: one config file, \texttt{llm\_config.py}, holds the model name.
\item Add what you built, with your hardware and your numbers, to the awesome list.
\end{enumerate}
\end{frame}
